\chapter{RECURSIVE-DESCENT PARSING AND THE ABSTRACT SYNTAX TREE}

This chapter presents the implementation of the recursive-descent parser and
the associated abstract syntax tree (AST) model. Starting from lexer-produced
tokens, the parser constructs a hierarchical \texttt{Program} AST that
represents expressions, statements, and control-flow structure in a form
suitable for downstream semantic analysis and execution. Implemented syntax
coverage includes arithmetic expressions, assignments, if-then-else, while,
function definitions, function calls, return, and built-in
\texttt{print(...)} statements. Equality operators (\texttt{==},
\texttt{!=}) are parsed as ordinary lowest-precedence expressions, so Boolean
results can appear in conditions, assignments, function arguments, and
\texttt{print(...)} expressions.

\section{Abstract syntax tree representation}
The AST is defined as a collection of Python dataclasses representing the
structural elements of the programming language. It serves as an intermediate
representation between parsing and later stages such as type checking and
interpretation.

A key distinction exists between tokens and AST nodes. Tokens are flat and
describe individual lexical elements such as operators or identifiers. In
contrast, the AST is hierarchical. For example, while a token may represent the
symbol \texttt{+}, a \texttt{BinaryOp} node represents an entire addition
operation, including its left and right operands as subtrees.

All AST nodes inherit from a common base class, allowing them to store
positional metadata (\texttt{line}, \texttt{column}). This ensures that errors
detected in later stages can still reference precise locations in the original
source code.

\begin{lstlisting}
@dataclass
class ASTNode:
   """Base class for every AST node. Carries source location for diagnostics."""
   line: int = field(default=0, kw_only=True)
   column: int = field(default=0, kw_only=True)
\end{lstlisting}

\subsection{Expressions}
Expressions represent computations or values within the program. A
\texttt{Literal} node stores values known at parse time, including integers,
floating-point numbers, booleans, and strings. String literals are stored
without surrounding quotation marks, as these are removed during parsing.

\begin{lstlisting}
if tok.token_type == TokenType.STRING_LITERAL:
   self._advance()
   return Literal(value=tok.lexeme[1:-1], line=tok.line, column=tok.column)
\end{lstlisting}

An \texttt{Identifier} node represents the use of a variable name within an
expression, indicating evaluation rather than declaration. A \texttt{BinaryOp}
node models arithmetic and comparison operations such as addition, subtraction,
multiplication, division, and equality checks, storing the operator along with
references to left and right operands. A \texttt{FunctionCall} node represents
function invocation, consisting of the function name and a list of argument
expressions.

\begin{lstlisting}
@dataclass
class BinaryOp(ASTNode):
   """An arithmetic or comparison binary operation.

   op is one of: '+' '-' '*' '/' '+.' '-.' '*.' '/.' '==' '!='
   Example:  a + b   x == 0   a +. b
   """
   op: str
   left: ASTNode
   right: ASTNode
\end{lstlisting}

\subsection{Statements}
Statements represent actions or control flow within the program. An
\texttt{Assign} node models variable assignment, pairing a variable name with a
value expression, while \texttt{Print} and \texttt{Return} nodes each store a
single expression and represent output and function return operations
respectively. A \texttt{Block} node groups multiple statements within braces,
representing a scoped execution sequence.

Control flow is handled via \texttt{If} and \texttt{While} nodes. The
\texttt{If} node stores a condition, a then-block, and an optional else-block,
while \texttt{While} stores a loop condition and body block. Function
definitions are represented by \texttt{FunctionDef} nodes, which include the
function name, a list of parameter names, and a body block. Parameter types
are intentionally not assigned at parse time, as type inference is handled
later by the type checker. Finally, \texttt{Program} serves as the AST root,
containing all top-level statements.

\section{Parser implementation}
The parser consumes the list of tokens generated by the lexer and produces a
single \texttt{Program} AST node. It is implemented as a hand-written
recursive-descent parser based on the project grammar. This avoids parser
generators and instead relies on explicit functions for each grammar rule.

The parser maintains internal state consisting of the token list and a position
index (\texttt{self.\_pos}) indicating the current token. Tokens are accessed
via helper methods such as \texttt{\_peek} and \texttt{\_advance}. Parsing
errors are reported using a custom \texttt{ParseError} exception with line and
column information consistent with lexer diagnostics.

\begin{lstlisting}
def parse(self) -> Program:
    """Parse the full token stream and return the root Program node."""
    line, col = self._peek().line, self._peek().column
    statements: list[ASTNode] = []
    while not self._is_at_end():
        statements.append(self._statement())
    return Program(statements=statements, line=line, column=col)
\end{lstlisting}

In addition to the main \texttt{parse()} entry point, the parser's statement
dispatcher (\texttt{\_statement}) implements a deterministic decision order
over token patterns: function definition, conditional, loop, return, print,
assignment, then expression statement fallback. This design is important
because it prevents ambiguity between syntactically similar forms. In
particular, assignments are recognized using one-token lookahead
(\texttt{IDENTIFIER} followed by \texttt{=}), while identifier-led expressions
that are not assignments are parsed as expression statements and must still end
with a semicolon. This gives the parser full coverage of both declarative-style
and expression-driven top-level statements without requiring a parser generator.

\subsection{Statement parsing}
The main parsing loop repeatedly processes statements until \texttt{EOF}, with
the statement type determined via dispatch on the current token. Function
definitions, conditionals, loops, returns, and prints each have dedicated
parsing methods. Assignments are detected using one-token lookahead
(\texttt{IDENTIFIER} followed by \texttt{ASSIGN}), preventing misclassification
of function calls as assignments. If no specific statement pattern matches, the
parser falls back to an expression statement and enforces a terminating
semicolon.

\begin{lstlisting}
# assignment: IDENTIFIER "=" expr ";"
if (
    tok.token_type == TokenType.IDENTIFIER
    and self._peek_next().token_type == TokenType.ASSIGN
):
    return self._assignment()
# fallback: bare expression statement  expr ";"
node = self._expr()
self._expect(TokenType.SEMICOLON, "Expected ';' after expression")
\end{lstlisting}

\subsection{Expression parsing and precedence}
Expressions are parsed in layered precedence levels: \texttt{\_expr} handles
comparison operators first (\texttt{==}, \texttt{!=}) after parsing an additive
left-hand side, then \texttt{\_additive} handles \texttt{+}, \texttt{-},
\texttt{+.}, and \texttt{-.}; \texttt{\_term} handles \texttt{*}, \texttt{/},
\texttt{*.}, and \texttt{/.}; and \texttt{\_factor} handles literals,
identifiers and function calls, and parenthesized expressions. Function calls
are recognised when an identifier is followed by \texttt{(}, with argument
lists parsed as comma-separated expressions. Parenthesized expressions are
supported to override precedence. Unary negation (\texttt{-x}, \texttt{-5}) is
also handled at the \texttt{\_factor} level: a leading \texttt{-} token causes
the parser to recursively parse the operand and return \texttt{BinaryOp(0 -
operand)}, so the existing type rules for subtraction cover negation without
any additional inference code. Float unary negation is handled similarly for
\texttt{-.}, producing \texttt{BinaryOp(0.0 -. operand)}.

\begin{lstlisting}
def _expr(self) -> ASTNode:      # ==, !=
   ...
def _additive(self) -> ASTNode:  # +, -, +., -.
   ...
def _term(self) -> ASTNode:  # *, /, *., /.
   ...
def _factor(self) -> ASTNode:  # unary -, literals, identifiers/calls, grouped expr
   ...
\end{lstlisting}

Figure~\ref{fig:ast-example} shows the AST produced for the statement
\texttt{z = x + y * 2;}. Because \texttt{\_term} handles \texttt{*} before
\texttt{\_additive} handles \texttt{+}, the multiplication sub-tree is nested one
level deeper, encoding the higher precedence of \texttt{*} over \texttt{+}
without any explicit priority table at runtime.

\begin{figure}[H]
\centering
\begin{tikzpicture}[
  anode/.style = {rectangle, rounded corners=3pt,
                  draw=nodeblue, fill=boxblue,
                  minimum width=2cm, minimum height=0.7cm,
                  align=center, font=\small\bfseries},
  leaf/.style  = {rectangle, rounded corners=3pt,
                  draw=arrowgray, fill=boxgray,
                  minimum width=1.8cm, minimum height=0.65cm,
                  align=center, font=\small\itshape},
  edge from parent/.style = {draw, -{Stealth[length=4pt]}, color=arrowgray},
  level 1/.style = {sibling distance=6cm, level distance=1.4cm},
  level 2/.style = {sibling distance=3.2cm, level distance=1.3cm},
  level 3/.style = {sibling distance=2.4cm, level distance=1.2cm},
]
\node[anode] {Assign (z)}
  child { node[anode] {BinaryOp (\texttt{+})}
    child { node[leaf] {Identifier (x)} }
    child { node[anode] {BinaryOp (\texttt{*})}
      child { node[leaf] {Identifier (y)} }
      child { node[leaf] {Literal (2)} }
    }
  };
\end{tikzpicture}
\caption{AST for \texttt{z = x + y * 2;}. Multiplication is a deeper sub-tree
than addition, encoding higher precedence through tree structure rather than
a runtime priority table.}
\label{fig:ast-example}
\end{figure}

This structural encoding of precedence eliminates the need for runtime
operator-priority resolution and keeps semantic analysis independent of
parsing strategy.

Expression parsing is intentionally split into layered methods to enforce
precedence and associativity in a transparent way. Parenthesized expressions
are parsed explicitly and re-enter \texttt{\_expr}, allowing precedence
override where required by source syntax. Comparison operators (\texttt{==},
\texttt{!=}) are parsed at the top of the \texttt{\_expr} method as the
lowest-precedence binary operators. Because comparisons are part of the
general expression grammar, Boolean results can appear in any expression
context --- as conditions, assignments, arguments, or print expressions.
Semantic enforcement that a condition is truly Boolean is deferred to the
type checker.

Figure~\ref{fig:ast-bool-example} shows the AST for an \texttt{if} statement
with a comparison condition, illustrating how the parser structures control flow
with a Boolean sub-tree.

\begin{figure}[H]
\centering
\begin{tikzpicture}[
  anode/.style = {rectangle, rounded corners=3pt,
                  draw=nodeblue, fill=boxblue,
                  minimum width=2cm, minimum height=0.7cm,
                  align=center, font=\small\bfseries},
  leaf/.style  = {rectangle, rounded corners=3pt,
                  draw=arrowgray, fill=boxgray,
                  minimum width=1.8cm, minimum height=0.65cm,
                  align=center, font=\small\itshape},
  edge from parent/.style = {draw, -{Stealth[length=4pt]}, color=arrowgray},
  level 1/.style = {sibling distance=5cm, level distance=1.4cm},
  level 2/.style = {sibling distance=2.8cm, level distance=1.3cm},
  level 3/.style = {sibling distance=2.4cm, level distance=1.2cm},
]
\node[anode] {If}
  child { node[anode] {BinaryOp (\texttt{==})}
    child { node[leaf] {Identifier (x)} }
    child { node[leaf] {Identifier (y)} }
  }
  child { node[anode] {Block (then)}
    child { node[anode] {Print}
      child { node[leaf] {Identifier (x)} }
    }
  };
\end{tikzpicture}
\caption{AST for \texttt{if (x == y) \{ print(x); \}}. The comparison produces
a \texttt{BinaryOp} node whose inferred type is \texttt{Boolean}, checked by
the type checker before execution.}
\label{fig:ast-bool-example}
\end{figure}

\subsection{Helper functions and lookahead}
As described above, the parser uses \texttt{\_check}, \texttt{\_expect},
\texttt{\_peek}, \texttt{\_peek\_next}, \texttt{\_advance}, and
\texttt{\_is\_at\_end} for token navigation and validation. Lookahead is central
for ambiguity resolution: assignment versus function call is resolved by
inspecting the next token, while comparison operators are recognised in
\texttt{\_expr} after parsing the left additive expression. Parser errors
follow a consistent diagnostic format:
\texttt{[line X, col Y] ParseError: <message>}.

The parser also includes dedicated list parsers for formal parameters
(\texttt{\_params}) and call arguments (\texttt{\_args}), both implemented as
comma-separated sequences with optional emptiness. This ensures consistent
handling of signatures such as \texttt{def f()} and \texttt{def f(a, b)} as
well as invocations like \texttt{f()} and \texttt{f(x, y+1)}.

\section{AST printing}
For debugging and reporting, an AST printer converts the tree into a readable
indented string, integrated into the pipeline output. The printer uses
visitor-style dynamic dispatch, mapping node class names to methods such as
\texttt{\_visit\_If} and \texttt{\_visit\_BinaryOp}, with output built as
multi-line text where indentation reflects tree depth. Node formatting
emphasises structure: \texttt{Program} and \texttt{Block} include statement
counts, \texttt{BinaryOp} explicitly shows left and right subtrees, and
\texttt{FunctionCall} prints each argument by index. \texttt{Literal} includes
both value and runtime Python type names for clarity.

\section{Parser and type checker boundary}
A deliberate architectural boundary is maintained between syntax construction
and semantic validation. The parser is responsible for producing a structurally
valid AST with accurate source locations, while type correctness and semantic
constraints are checked later by the type checker. For example,
\texttt{\_expr} may syntactically produce a node whose eventual type is not
Boolean, but rejection of non-Boolean control-flow conditions is handled during
semantic analysis. This separation keeps parser logic focused on grammar
conformance and allows a clear responsibility split across team member
contributions.
